\subsection{Token Accounting and Replay Capacity}
Table~\ref{tab:app-ablation-tokens} uses the first run of each configuration; it is separate from the three-run means above. We sum recorded input and output tokens, including retries and invalid completions with known usage. Exact-request cache hits incur no additional tokens. The totals cover run-local initialization, initial assessment, optimization, and final assessment when present. All component configurations use the same initial predicates. The totals reflect recorded usage rather than imputed costs. The iteration-only column excludes initialization and initial/final assessment. Attempts with unavailable token use are not assigned zero consumption: they are excluded from known totals and counted separately.

\begin{table*}[t]
\centering
\small
\begin{tabular}{lrrrrr}
\toprule
Variant & F1 & nDCG@1 & Total (M) & Iterations (M) & Unknown \\
\midrule
Both Disabled & 0.379 & 0.553 & 840.052 & 662.211 & 12 \\
Only Gate & 0.494 & 0.643 & 866.352 & 701.433 & 13 \\
Only MBR & \underline{0.581} & \underline{0.684} & 1181.026 & 1011.920 & 20 \\
Both Enabled & \textbf{0.630} & \textbf{0.742} & 1324.251 & 1073.904 & 16 \\
\bottomrule
\end{tabular}
\caption{Single-run final-test quality and known token use for the component configurations in Figure~\ref{fig:performance_token_usage}. Bold and underline mark the best and second-best quality values. M denotes one million tokens. Total includes all recorded run-local phases; Iterations includes positive-numbered optimization iterations only. Unknown counts attempts whose token use is unavailable.}
\label{tab:app-ablation-tokens}
\end{table*}

We additionally vary the per-class replay capacity over 1, 3, 5, and 10 while retaining MBR and the held-out gate (Table~\ref{tab:app-replay-capacity}). The replay set can contain up to twice this capacity plus the current task if it is not already retained. Larger capacity increases known token use. The largest tested capacity has the highest final F1 and nDCG@10, whereas capacity five has the highest nDCG@1. The sensitivity results therefore depend on which aspect of routing quality is measured.

\begin{table*}[t]
\centering
\small
\begin{tabular}{lrrrrr}
\toprule
Capacity & F1 & nDCG@1 & nDCG@10 & Tokens (M) & Unknown \\
\midrule
1 & 0.534 & 0.659 & 0.739 & 949.061 & 28 \\
3 & 0.579 & 0.702 & 0.769 & 1175.955 & 27 \\
5 & \underline{0.630} & \textbf{0.742} & \underline{0.800} & 1324.251 & 16 \\
10 & \textbf{0.641} & \underline{0.730} & \textbf{0.805} & 1809.623 & 46 \\
\bottomrule
\end{tabular}
\caption{Replay-capacity sensitivity using one run per capacity, Qwen3.5 4B, SkillSight, and both components enabled. Capacity is per relevance class. Quality is measured on the final policy over 748 test tasks. Bold and underline mark the best and second-best quality values.}
\label{tab:app-replay-capacity}
\end{table*}
